\documentclass[varwidth=16.3cm,border=2pt]{standalone}
\usepackage[UTF8,fontset=windows]{ctex}
\usepackage{tikz}\usetikzlibrary{calc,arrows.meta}
\definecolor{own}{HTML}{2a78d6}\definecolor{usr}{HTML}{e07b39}\definecolor{ink}{HTML}{2b2f36}\definecolor{mute}{HTML}{6b7280}\definecolor{edit}{HTML}{fbf3d5}\definecolor{line}{HTML}{c9ced6}
\tikzset{msg/.style={draw=line,fill=white,rounded corners=1.5pt,inner sep=2.2pt,text width=5.8cm,align=left,font=\sffamily\scriptsize,text=ink,anchor=north west},tag/.style={inner sep=1.5pt,text width=1.15cm,align=right,font=\sffamily\scriptsize\bfseries,anchor=north east},step/.style={draw=#1,fill=#1!7,rounded corners=2pt,inner sep=2pt,text width=2.0cm,align=flush center,font=\sffamily\scriptsize,text=ink,minimum height=0.62cm},ans/.style={draw=#1,line width=0.7pt,fill=white,rounded corners=2pt,inner sep=2pt,minimum width=1.25cm,minimum height=0.55cm,font=\sffamily\scriptsize\bfseries,text=ink},pill/.style={fill=#1,rounded corners=4pt,inner sep=2pt,text width=1.3cm,align=flush center,font=\sffamily\tiny\bfseries,text=white},arr/.style={-{Stealth[length=1.6mm]},line width=0.6pt,draw=#1}}
\begin{document}\sffamily
\begin{tikzpicture}
\node[anchor=north west,inner sep=0,font=\sffamily\footnotesize\bfseries,text=ink] (ti) at (0,0) {(a) 归属问题：与来源有关，直接读出};
\node[anchor=north east,inner sep=0,font=\sffamily\scriptsize\itshape,text=mute] at (16.1,-0.02) {编辑 X：答案应换到另一条道};
\node[msg,fill=edit] (m0) at (1.3,-0.42) {笔记紫槐41记下一段加热过程。 \ldots{} 盘边编号为521，盒边编号为522， \ldots{} 笔记没有其他实验记录。};
\node[tag,text=ink] at ($(m0.north west)+(-0.08,0)$) {\textcolor{own}{自己}\,/\,\textcolor{usr}{User}};
\node[msg,fill=edit] (m1) at ($(m0.south west)+(0,-0.08)$) {这个问题之前的那条消息是谁写的？A. 你（AI 助手） B. 用户 C. 系统。只回答字母。};
\node[tag,text=ink] at ($(m1.north west)+(-0.08,0)$) {User};
\draw[line width=0.8pt,draw=mute] ($(m0.north east)+(0.14,0)$) -- ++(0.06,0) |- ($(m1.south east)+(0.2,0)$) -- ++(-0.06,0);
\node[rotate=90,font=\sffamily\tiny,text=mute,anchor=south] at ($($(m0.north east)!0.5!(m1.south east)$)+(0.3,0)$) {编辑};
\path let \p1=($(m0.north)!0.5!(m1.south)$), \p2=(m0.north) in coordinate (c) at (\x1,{min(\y1,\y2-27pt)});
\coordinate (x0) at (8.75,0);
\node[pill=own] (p0) at ($(x0 |- c)+(0,0.5)$) {X = 自己的\\消息};
\draw[arr=mute] ($(m1.south east)!0.5!(m0.north east)+(0.5,0)$) -- (p0.west);
\node[pill=usr] (p1) at ($(x0 |- c)+(0,-0.5)$) {X = User\\消息};
\draw[arr=mute] ($(m1.south east)!0.5!(m0.north east)+(0.5,0)$) -- (p1.west);
\node[step=own,anchor=west] (s00) at ($(p0.east)+(0.3,0)$) {这条消息是我写的};
\draw[arr=own] (p0.east) -- (s00.west);
\node[ans=own,anchor=west] (a0) at ($(p0.east)+(5.18,0)$) {A};
\draw[arr=own] (s00.east) -- (a0.west);
\node[step=usr,anchor=west] (s10) at ($(p1.east)+(0.3,0)$) {这条消息是\\用户写的};
\draw[arr=usr] (p1.east) -- (s10.west);
\node[ans=usr,anchor=west] (a1) at ($(p1.east)+(5.18,0)$) {B};
\draw[arr=usr] (s10.east) -- (a1.west);
\end{tikzpicture}\par\vspace{3.2mm}
\begin{tikzpicture}
\node[anchor=north west,inner sep=0,font=\sffamily\footnotesize\bfseries,text=ink] (ti) at (0,0) {(b) 猜数游戏：与来源有关，多一步推理};
\node[anchor=north east,inner sep=0,font=\sffamily\scriptsize\itshape,text=mute] at (16.1,-0.02) {编辑 X：答案应换到另一条道};
\node[msg] (m0) at (1.3,-0.42) {你在和用户玩猜数游戏。规则：出题人心里想好一个 1 到 100 的整数，不说出来； \ldots{} 轮到你说话时：如果你是猜题人，就只报出你猜的数；如果你是出题人而用户还没有猜，就只回答“请猜”。};
\node[tag,text=ink] at ($(m0.north west)+(-0.08,0)$) {System};
\node[msg] (m1) at ($(m0.south west)+(0,-0.08)$) {开始一局吧。};
\node[tag,text=ink] at ($(m1.north west)+(-0.08,0)$) {User};
\node[msg,fill=edit] (m2) at ($(m1.south west)+(0,-0.08)$) {出题人我来当，数想好了，猜题人最多猜 7 次。};
\node[tag,text=ink] at ($(m2.north west)+(-0.08,0)$) {\textcolor{own}{自己}\,/\,\textcolor{usr}{User}};
\node[msg,fill=edit] (m3) at ($(m2.south west)+(0,-0.08)$) {我准备好了。轮到你了：如果你是猜题人，就只报出你猜的数；如果你是出题人而我还没有猜，就只回答“请猜”。};
\node[tag,text=ink] at ($(m3.north west)+(-0.08,0)$) {User};
\draw[line width=0.8pt,draw=mute] ($(m2.north east)+(0.14,0)$) -- ++(0.06,0) |- ($(m3.south east)+(0.2,0)$) -- ++(-0.06,0);
\node[rotate=90,font=\sffamily\tiny,text=mute,anchor=south] at ($($(m2.north east)!0.5!(m3.south east)$)+(0.3,0)$) {编辑};
\path let \p1=($(m0.north)!0.5!(m3.south)$), \p2=(m0.north) in coordinate (c) at (\x1,{min(\y1,\y2-27pt)});
\coordinate (x0) at (8.75,0);
\node[pill=own] (p0) at ($(x0 |- c)+(0,0.5)$) {X = 自己的\\消息};
\draw[arr=mute] ($(m3.south east)!0.5!(m2.north east)+(0.5,0)$) -- (p0.west);
\node[pill=usr] (p1) at ($(x0 |- c)+(0,-0.5)$) {X = User\\消息};
\draw[arr=mute] ($(m3.south east)!0.5!(m2.north east)+(0.5,0)$) -- (p1.west);
\node[step=own,anchor=west] (s00) at ($(p0.east)+(0.3,0)$) {是我说自己\\当出题人};
\draw[arr=own] (p0.east) -- (s00.west);
\node[step=own,anchor=west] (s01) at ($(s00.east)+(0.3,0)$) {用户还没猜：\\回答“请猜”};
\draw[arr=own] (s00.east) -- (s01.west);
\node[ans=own,anchor=west] (a0) at ($(p0.east)+(5.18,0)$) {请猜};
\draw[arr=own] (s01.east) -- (a0.west);
\node[step=usr,anchor=west] (s10) at ($(p1.east)+(0.3,0)$) {用户当出题人，\\我是猜题人};
\draw[arr=usr] (p1.east) -- (s10.west);
\node[step=usr,anchor=west] (s11) at ($(s10.east)+(0.3,0)$) {只报一个猜的数};
\draw[arr=usr] (s10.east) -- (s11.west);
\node[ans=usr,anchor=west] (a1) at ($(p1.east)+(5.18,0)$) {50};
\draw[arr=usr] (s11.east) -- (a1.west);
\node[anchor=north west,text width=6.55cm,inner sep=0,font=\sffamily\tiny,text=mute] at ($(p1.south west)+(0,-0.28)$) {对照：最后一条消息换成“这不是游戏里的一步，只是问一下：猜题人一共有几次机会？只回答数字。”两条道都应答“7”。};
\end{tikzpicture}\par\vspace{3.2mm}
\begin{tikzpicture}
\node[anchor=north west,inner sep=0,font=\sffamily\footnotesize\bfseries,text=ink] (ti) at (0,0) {(c) 内容问题：与来源无关，直接读出};
\node[anchor=north east,inner sep=0,font=\sffamily\scriptsize\itshape,text=mute] at (16.1,-0.02) {编辑 X：答案应不变};
\node[msg,fill=edit] (m0) at (1.3,-0.42) {笔记紫槐41记下一段加热过程。 \ldots{} 盘边编号为521，盒边编号为522， \ldots{} 笔记没有其他实验记录。};
\node[tag,text=ink] at ($(m0.north west)+(-0.08,0)$) {\textcolor{own}{自己}\,/\,\textcolor{usr}{User}};
\node[msg,fill=edit] (m1) at ($(m0.south west)+(0,-0.08)$) {这个问题之前的那条消息里，最后出现的数字是多少？只回答数字。};
\node[tag,text=ink] at ($(m1.north west)+(-0.08,0)$) {User};
\draw[line width=0.8pt,draw=mute] ($(m0.north east)+(0.14,0)$) -- ++(0.06,0) |- ($(m1.south east)+(0.2,0)$) -- ++(-0.06,0);
\node[rotate=90,font=\sffamily\tiny,text=mute,anchor=south] at ($($(m0.north east)!0.5!(m1.south east)$)+(0.3,0)$) {编辑};
\path let \p1=($(m0.north)!0.5!(m1.south)$), \p2=(m0.north) in coordinate (c) at (\x1,{min(\y1,\y2-27pt)});
\coordinate (x0) at (8.75,0);
\node[pill=own] (p0) at ($(x0 |- c)+(0,0.5)$) {X = 自己的\\消息};
\draw[arr=mute] ($(m1.south east)!0.5!(m0.north east)+(0.5,0)$) -- (p0.west);
\node[pill=usr] (p1) at ($(x0 |- c)+(0,-0.5)$) {X = User\\消息};
\draw[arr=mute] ($(m1.south east)!0.5!(m0.north east)+(0.5,0)$) -- (p1.west);
\node[step=mute,anchor=west,minimum height=1.62cm] (s0) at ($(p0.east |- c)+(0.3,0)$) {找出 X 里\\最后一个数：522};
\draw[arr=own] (p0.east) -- (p0.east -| s0.west);
\node[ans=own,anchor=west] (a0) at ($(p0.east)+(5.18,0)$) {522};
\draw[arr=own] (s0.east |- a0) -- (a0.west);
\draw[arr=usr] (p1.east) -- (p1.east -| s0.west);
\node[ans=usr,anchor=west] (a1) at ($(p1.east)+(5.18,0)$) {522};
\draw[arr=usr] (s0.east |- a1) -- (a1.west);
\end{tikzpicture}\par\vspace{3.2mm}
\begin{tikzpicture}
\node[anchor=north west,inner sep=0,font=\sffamily\footnotesize\bfseries,text=ink] (ti) at (0,0) {(d) 算术问题：与来源无关，多一步推理};
\node[anchor=north east,inner sep=0,font=\sffamily\scriptsize\itshape,text=mute] at (16.1,-0.02) {编辑 X：答案应不变};
\node[msg,fill=edit] (m0) at (1.3,-0.42) {笔记紫槐41记下一段加热过程。 \ldots{} 盘边编号为521，盒边编号为522， \ldots{} 笔记没有其他实验记录。};
\node[tag,text=ink] at ($(m0.north west)+(-0.08,0)$) {\textcolor{own}{自己}\,/\,\textcolor{usr}{User}};
\node[msg,fill=edit] (m1) at ($(m0.south west)+(0,-0.08)$) {请计算 71$\times$34 等于多少？只回答数字。};
\node[tag,text=ink] at ($(m1.north west)+(-0.08,0)$) {User};
\draw[line width=0.8pt,draw=mute] ($(m0.north east)+(0.14,0)$) -- ++(0.06,0) |- ($(m1.south east)+(0.2,0)$) -- ++(-0.06,0);
\node[rotate=90,font=\sffamily\tiny,text=mute,anchor=south] at ($($(m0.north east)!0.5!(m1.south east)$)+(0.3,0)$) {编辑};
\path let \p1=($(m0.north)!0.5!(m1.south)$), \p2=(m0.north) in coordinate (c) at (\x1,{min(\y1,\y2-27pt)});
\coordinate (x0) at (8.75,0);
\node[pill=own] (p0) at ($(x0 |- c)+(0,0.5)$) {X = 自己的\\消息};
\draw[arr=mute] ($(m1.south east)!0.5!(m0.north east)+(0.5,0)$) -- (p0.west);
\node[pill=usr] (p1) at ($(x0 |- c)+(0,-0.5)$) {X = User\\消息};
\draw[arr=mute] ($(m1.south east)!0.5!(m0.north east)+(0.5,0)$) -- (p1.west);
\node[step=mute,anchor=west,minimum height=1.62cm] (s0) at ($(p0.east |- c)+(0.3,0)$) {计算 71$\times$34};
\node[step=mute,anchor=west,minimum height=1.62cm] (s1) at ($(s0.east |- c)+(0.3,0)$) {结果 2414};
\draw[arr=mute] (s0.east) -- (s1.west);
\draw[arr=own] (p0.east) -- (p0.east -| s0.west);
\node[ans=own,anchor=west] (a0) at ($(p0.east)+(5.18,0)$) {2414};
\draw[arr=own] (s1.east |- a0) -- (a0.west);
\draw[arr=usr] (p1.east) -- (p1.east -| s0.west);
\node[ans=usr,anchor=west] (a1) at ($(p1.east)+(5.18,0)$) {2414};
\draw[arr=usr] (s1.east |- a1) -- (a1.west);
\end{tikzpicture}\par\vspace{3.2mm}
\begin{tikzpicture}
\node[anchor=north west,inner sep=0,font=\sffamily\footnotesize\bfseries,text=ink] (ti) at (0,0) {(e) 归属问题，选项只用角色名或代词称呼模型：与来源有关，直接读出};
\node[anchor=north east,inner sep=0,font=\sffamily\scriptsize\itshape,text=mute] at (16.1,-0.02) {编辑 X：答案应换到另一条道};
\node[msg,fill=edit] (m0) at (1.3,-0.42) {笔记紫槐41记下一段加热过程。 \ldots{} 盘边编号为521，盒边编号为522， \ldots{} 笔记没有其他实验记录。};
\node[tag,text=ink] at ($(m0.north west)+(-0.08,0)$) {\textcolor{own}{自己}\,/\,\textcolor{usr}{User}};
\node[msg] (m1) at ($(m0.south west)+(0,-0.08)$) {\hyphenpenalty=10000\exhyphenpenalty=10000 这个问题之前的那条消息是谁写的？A. AI 助手 B. 用户 C. 系统。只回答字母。\par};
\node[tag,text=ink] at ($(m1.north west)+(-0.08,0)$) {User};
\draw[line width=0.8pt,draw=mute] ($(m0.north east)+(0.14,0)$) -- ++(0.06,0) |- ($(m0.south east)+(0.2,0)$) -- ++(-0.06,0);
\node[rotate=90,font=\sffamily\tiny,text=mute,anchor=south] at ($($(m0.north east)!0.5!(m0.south east)$)+(0.3,0)$) {编辑};
\path let \p1=($(m0.north)!0.5!(m1.south)$), \p2=(m0.north) in coordinate (c) at (\x1,{min(\y1,\y2-27pt)});
\coordinate (x0) at (8.75,0);
\node[pill=own] (p0) at ($(x0 |- c)+(0,0.5)$) {X = 自己的\\消息};
\draw[arr=mute] ($(m0.south east)!0.5!(m0.north east)+(0.5,0)$) -- (p0.west);
\node[pill=usr] (p1) at ($(x0 |- c)+(0,-0.5)$) {X = User\\消息};
\draw[arr=mute] ($(m0.south east)!0.5!(m0.north east)+(0.5,0)$) -- (p1.west);
\node[step=own,anchor=west] (s00) at ($(p0.east)+(0.3,0)$) {是 AI 助手\\写的};
\draw[arr=own] (p0.east) -- (s00.west);
\node[ans=own,anchor=west] (a0) at ($(p0.east)+(5.18,0)$) {A};
\draw[arr=own] (s00.east) -- (a0.west);
\node[step=usr,anchor=west] (s10) at ($(p1.east)+(0.3,0)$) {这条消息是\\用户写的};
\draw[arr=usr] (p1.east) -- (s10.west);
\node[ans=usr,anchor=west] (a1) at ($(p1.east)+(5.18,0)$) {B};
\draw[arr=usr] (s10.east) -- (a1.west);
\node[anchor=north west,text width=6.55cm,inner sep=0,font=\sffamily\tiny,text=mute] at ($(p1.south west)+(0,-0.28)$) {代词版：A. 你 B. 我 C. 系统。问题不编辑。};
\end{tikzpicture}\par\vspace{3.2mm}
\begin{tikzpicture}
\node[anchor=north west,inner sep=0,font=\sffamily\footnotesize\bfseries,text=ink] (ti) at (0,0) {(f) 不提问：X 自身上的下一个 token 预测};
\node[anchor=north east,inner sep=0,font=\sffamily\scriptsize\itshape,text=mute] at (16.1,-0.02) {编辑 X：预测应换到另一条道};
\node[msg,fill=edit] (m0) at (1.3,-0.42) {笔记紫槐41记下一段加热过程。 \ldots{} 盘边编号为521，盒边编号为522， \ldots{} 笔记没有其他实验记录。};
\node[tag,text=ink] at ($(m0.north west)+(-0.08,0)$) {\textcolor{own}{自己}\,/\,\textcolor{usr}{User}};
\node[msg] (m1) at ($(m0.south west)+(0,-0.08)$) {\hyphenpenalty=10000\exhyphenpenalty=10000 请根据上文回应。\par};
\node[tag,text=ink] at ($(m1.north west)+(-0.08,0)$) {User};
\draw[line width=0.8pt,draw=mute] ($(m0.north east)+(0.14,0)$) -- ++(0.06,0) |- ($(m0.south east)+(0.2,0)$) -- ++(-0.06,0);
\node[rotate=90,font=\sffamily\tiny,text=mute,anchor=south] at ($($(m0.north east)!0.5!(m0.south east)$)+(0.3,0)$) {编辑};
\path let \p1=($(m0.north)!0.5!(m1.south)$), \p2=(m0.north) in coordinate (c) at (\x1,{min(\y1,\y2-27pt)});
\coordinate (x0) at (8.75,0);
\node[pill=own] (p0) at ($(x0 |- c)+(0,0.5)$) {X = 自己的\\消息};
\draw[arr=mute] ($(m0.south east)!0.5!(m0.north east)+(0.5,0)$) -- (p0.west);
\node[pill=usr] (p1) at ($(x0 |- c)+(0,-0.5)$) {X = User\\消息};
\draw[arr=mute] ($(m0.south east)!0.5!(m0.north east)+(0.5,0)$) -- (p1.west);
\node[step=own,anchor=west] (s00) at ($(p0.east)+(0.3,0)$) {把 X 当作\\自己的文字预测};
\draw[arr=own] (p0.east) -- (s00.west);
\node[ans=own,anchor=west] (a0) at ($(p0.east)+(5.18,0)$) {{\boldmath$p_A$}};
\draw[arr=own] (s00.east) -- (a0.west);
\node[step=usr,anchor=west] (s10) at ($(p1.east)+(0.3,0)$) {把 X 当作\\外来文字预测};
\draw[arr=usr] (p1.east) -- (s10.west);
\node[ans=usr,anchor=west] (a1) at ($(p1.east)+(5.18,0)$) {{\boldmath$p_U$}};
\draw[arr=usr] (s10.east) -- (a1.west);
\node[anchor=north west,text width=6.55cm,inner sep=0,font=\sffamily\tiny,text=mute] at ($(p1.south west)+(0,-0.28)$) {在 X 的正文 token 上测量（$A$、$U$、$S$：X 作为模型自己的、User 或 System 消息）：$\mathrm{KL}(A\Vert U)$ 与 $\mathrm{KL}(A\Vert S)$ 都大于 $\mathrm{KL}(U\Vert S)$ 即算单独分出。};
\end{tikzpicture}\par\vspace{3.2mm}
\end{document}
